% Abhakseite „Das kann ich“ (erzeugt von verz.py)
\eng
\hypertarget{abhaken}{}\einheitenkopf{Das kann ich}
\begingroup\setlength{\parskip}{3pt}
\noindent Hake ab, was du kannst.\par\medskip
\par\medskip\noindent\textbf{Kennst du schon}\par\smallskip
\noindent\hangindent=2.2em\hangafter=1 $\square$\quad 1. Ich kann eine Strichliste auszählen und angeben, welcher Anteil das ist.\par
\noindent\hangindent=2.2em\hangafter=1 $\square$\quad 2. Ich kann ausrechnen, wie viel Prozent ein Teil vom Ganzen ist.\par
\noindent\hangindent=2.2em\hangafter=1 $\square$\quad 3. Ich kann Werte aus einem Säulendiagramm ablesen.\par
\noindent\hangindent=2.2em\hangafter=1 $\square$\quad 4. Ich kann Zahlen ordnen und Werte an einer Skala ablesen.\par
\noindent\hangindent=2.2em\hangafter=1 $\square$\quad 5. Ich kann Längen in Zentimetern und Millimetern angeben und zeichnen.\par
\noindent\hangindent=2.2em\hangafter=1 $\square$\quad 6. Ich kann Winkel als Teil des Vollkreises sehen und mit dem Geodreieck zeichnen.\par
\noindent\hangindent=2.2em\hangafter=1 $\square$\quad 7. Ich kann Winkel als Teil des Vollkreises sehen und mit dem Geodreieck zeichnen – weiter.\par
\noindent\hangindent=2.2em\hangafter=1 $\square$\quad 8. Ich kann mit Dezimalzahlen rechnen und sinnvoll runden.\par
\noindent\hangindent=2.2em\hangafter=1 $\square$\quad 9. Ich kann den kleinsten und den größten Wert, die Spannweite, den Mittelwert und den Median einer Liste bestimmen.\par
\noindent\hangindent=2.2em\hangafter=1 $\square$\quad 10. Ich kann ausrechnen, was „die Hälfte“, „ein Drittel mehr“ oder „doppelt so viel“ bedeutet.\par
\noindent\hangindent=2.2em\hangafter=1 $\square$\quad 11. Ich finde den Fehler, wenn ein Anteil vom falschen Ganzen berechnet wurde.\par
\par\medskip\noindent\textbf{Einheit 1 · Streifen- und Kreisdiagramm}\par\smallskip
\noindent\hangindent=2.2em\hangafter=1 $\square$\quad 12. Ich erkenne, was das Ganze ist.\par
\noindent\hangindent=2.2em\hangafter=1 $\square$\quad 13. Ich erkenne, ob eine Zahl ein Anteil in Prozent oder ein Winkel in Grad ist.\par
\noindent\hangindent=2.2em\hangafter=1 $\square$\quad 14. Ich kann einen Anteil als Streifen darstellen.\par
\noindent\hangindent=2.2em\hangafter=1 $\square$\quad 15. Ich kann einen Anteil als Streifen darstellen – weiter.\par
\noindent\hangindent=2.2em\hangafter=1 $\square$\quad 16. Ich kann ausrechnen, welcher Winkel im Kreisdiagramm zu einem Anteil gehört.\par
\noindent\hangindent=2.2em\hangafter=1 $\square$\quad 17. Ich kann ein Kreisdiagramm zeichnen (kennst du wahrscheinlich seit Klasse 6).\par
\noindent\hangindent=2.2em\hangafter=1 $\square$\quad 18. Ich kann ein Kreisdiagramm zeichnen – weiter.\par
\noindent\hangindent=2.2em\hangafter=1 $\square$\quad 19. Ich kann schätzen, welcher Anteil zu einem Sektor gehört.\par
\noindent\hangindent=2.2em\hangafter=1 $\square$\quad 20. Ich kann die Sektoren eines Kreisdiagramms den Anteilen zuordnen.\par
\noindent\hangindent=2.2em\hangafter=1 $\square$\quad 21. Ich finde den Fehler beim Winkel im Kreisdiagramm.\par
\noindent\hangindent=2.2em\hangafter=1 $\square$\quad 22. Ich kann begründen, wie Anteil und Winkel im Kreisdiagramm zusammenhängen.\par
\noindent\hangindent=2.2em\hangafter=1 $\square$\quad 23. Ich kann entscheiden, ob ein Kreisdiagramm zu einer Umfrage passt.\par
\par\medskip\noindent\textbf{Einheit 2 · Diagramme beurteilen und Boxplot}\par\smallskip
\noindent\hangindent=2.2em\hangafter=1 $\square$\quad 24. Ich erkenne, wo die senkrechte Achse eines Diagramms anfängt.\par
\noindent\hangindent=2.2em\hangafter=1 $\square$\quad 25. Ich kann prüfen, ob eine Aussage zu einem Diagramm stimmt.\par
\noindent\hangindent=2.2em\hangafter=1 $\square$\quad 26. Ich kann prüfen, ob eine Aussage zu einem Diagramm stimmt – weiter.\par
\noindent\hangindent=2.2em\hangafter=1 $\square$\quad 27. Ich kann erklären, warum ein Diagramm täuscht.\par
\noindent\hangindent=2.2em\hangafter=1 $\square$\quad 28. Ich kann einen Trend in einem Diagramm beurteilen (kommt in Klasse 9, am Gymnasium neu in diesem Jahr).\par
\noindent\hangindent=2.2em\hangafter=1 $\square$\quad 29. Ich kann ein Diagramm mit einer Achse ab 0 zeichnen und eine Zeitungsaussage dazu prüfen.\par
\noindent\hangindent=2.2em\hangafter=1 $\square$\quad 30. Ich kann einen Boxplot lesen (kommt nächstes Jahr, je nach Buch bis Klasse 9; am Gymnasium neu in diesem Jahr).\par
\noindent\hangindent=2.2em\hangafter=1 $\square$\quad 31. Ich kann einen Boxplot zeichnen (kommt nächstes Jahr, je nach Buch bis Klasse 9; am Gymnasium neu in diesem Jahr).\par
\noindent\hangindent=2.2em\hangafter=1 $\square$\quad 32. Ich finde den Fehler, wenn jemand nur die Säulenhöhe vergleicht.\par
\noindent\hangindent=2.2em\hangafter=1 $\square$\quad 33. Ich kann begründen, wann die Achse eines Diagramms bei 0 beginnen sollte.\par
\noindent\hangindent=2.2em\hangafter=1 $\square$\quad 34. Ich kann mit Zahlen aus einer Zählung einen Antrag begründen.\par
\endgroup
